% moi_r2: I weights each mass by r^2 (Ch 14, fig:moir2)
% Three equal masses m at r, 2r, 3r from a rotation axis (r = 2.7 cm drawn).  Above
% each, a bar of height proportional to its contribution m r^2: 1 : 4 : 9
% (0.24 cm per unit).  Distances use the same dimension style as the moid_* figures,
% and the axis the same dash-dot line.
\begin{tikzpicture}[x=1cm,y=1cm]
  \def\u{2.7}    % drawn length of r
  \def\h{0.24}   % bar height per unit of m r^2
  \def\bw{0.26}  % half width of a bar
  \draw[rotation axis] (0,-1.85) -- (0,2.85) node[anchor=south] {axis};
  % masses, bars and labels
  \foreach \n in {1,2,3}{%
    \pgfmathsetmacro\x{\n*\u}
    \pgfmathsetmacro\top{0.45+\h*\n*\n}
    \draw[cPrim,thin line,fill=cPrim!25] (\x-\bw,0.45) rectangle (\x+\bw,\top);
    \pgfmathtruncatemacro\sq{\n*\n}
    \node[cPrim,anchor=south] at (\x,\top) {$\sq\times$};
    \draw[body] (\x,0) circle[radius=0.23];
    \node at (\x,0) {$m$};
  }
  % what the bars measure
  \node[cPrim,anchor=west] at (0.25,2.55) {contribution to $I$: $mr^2$};
  % distances from the axis, stacked
  \foreach \n/\y/\lab in {1/-0.55/r,2/-1.0/2r,3/-1.45/3r}{%
    \draw[dimtick] (\n*\u,-0.26) -- (\n*\u,\y-0.12);
    \draw[dim] (0,\y) -- (\n*\u,\y);
    \node[cSteel,anchor=west,inner sep=2pt] at (\n*\u,\y) {$\lab$};
  }
\end{tikzpicture}
